\chapter{Detailed Design}
\label{app:design}

This appendix records the implementation detail that Chapter~\ref{ch:method}
deliberately omits: the module map, the data schema every stage exchanges,
the backend abstraction that lets the same pipeline run on mocks or on a live
browser, and the exact configuration used for the runs reported in
Chapter~\ref{ch:experiments}.

\section{Module map}
\label{app:design-modules}

Table~\ref{tab:app-modules} lists every module of the implementation with its
responsibility. The pipeline core is deliberately small: the whole discovery
path from a raw trajectory to an admitted skill is under 700 lines, which is
what makes the deterministic mock tests practical.

{\footnotesize
\begin{longtable}{@{}L{4.3cm}rL{8.0cm}@{}}
\caption[Module map of the HSGE implementation]{Module map. Line counts are of
the implementation only, excluding tests.}
\label{tab:app-modules}\\
\toprule
\textbf{Module} & \textbf{Lines} & \textbf{Responsibility} \\
\midrule
\endfirsthead
\toprule
\textbf{Module} & \textbf{Lines} & \textbf{Responsibility} \\
\midrule
\endhead
\bottomrule
\endfoot
\multicolumn{3}{@{}l}{\emph{Pipeline core (\texttt{hsge/})}} \\
\midrule
\texttt{schema.py}      & 129 & Shared data types: \texttt{Step},
\texttt{SubTrajectory}, \texttt{SkillCard}, \texttt{ParamSlot}, \texttt{Task} \\
\texttt{segmentation.py} & 138 & Candidate cuts (URL change + accessibility-tree
similarity drop), LLM arbitration, single-step navigation merge \\
\texttt{clustering.py}  & 100 & Streaming nonparametric clusterer
(cosine-threshold / DP-means), no preset cluster count \\
\texttt{abstraction.py} &  59 & Lifts a mature cluster into a parameterized
\texttt{SkillCard} \\
\texttt{verification.py} &  62 & Grounds, executes and admits a skill only on
$\geq k$ successes in held-out contexts \\
\texttt{library.py}     & 105 & \texttt{SkillLibrary}: observe $\rightarrow$
cluster $\rightarrow$ abstract $\rightarrow$ verify $\rightarrow$ admit;
retrieval and pruning \\
\texttt{discovery.py}   &  90 & \texttt{Discovery.on\_trajectory}: the online
hook that drives the pipeline from inside any exploration loop \\
\texttt{graph.py}       &  61 & \texttt{HierarchicalGraph}: skill nodes and
skill$\rightarrow$skill composition edges \\
\texttt{operators.py}   & 128 & Precondition/effect operators, BFS planner,
\texttt{all\_chains} bounded DFS over the operator graph \\
\texttt{operator\_induction.py} & 140 & LLM induction of preconditions/effects
from a skill card plus boundary observations; grounding check \\
\texttt{predicate\_registry.py} &  66 & Canonicalizes free-form NL predicates
into stable ids by cosine threshold \\
\texttt{curriculum.py}  &  60 & Composes skills into easy$\rightarrow$hard tasks \\
\texttt{transfer.py}    &  53 & Retrieve, re-ground and verify a skill on a new
site \\
\texttt{llm.py}         & 149 & \texttt{LLM.complete\_json} and
\texttt{Embedder.encode} interfaces; \texttt{MockLLM},
\texttt{HashingEmbedder} \\
\texttt{backends.py}    & 128 & Real backends: \texttt{OpenAILLM} (retrying,
optionally vision-capable) and \texttt{STEmbedder} \\
\texttt{env.py}         &  95 & \texttt{WebEnv} interface
(\texttt{reset}/\texttt{step}/\texttt{check}) and the mock sites \\
\midrule
\multicolumn{3}{@{}l}{\emph{WebArena runtime (\texttt{hsge/webarena/})}} \\
\midrule
\texttt{runtime.py}     & 681 & \texttt{WebArenaExplorer}: frontier loop,
skill$\rightarrow$operator promotion, goal tasks, persistence \\
\texttt{agents.py}      & 363 & \texttt{PageExplorer} (task proposal) and
\texttt{Solver} (action loop), prompts, stuck detection \\
\texttt{browser.py}     & 143 & \texttt{WebArenaBrowser}: BrowserGym wrapper,
observations, screenshots, href mining \\
\texttt{auth.py}        & 138 & Per-site login flows and cached
\texttt{storage\_state} sessions \\
\texttt{predicates.py}  &  61 & URL/DOM terminal-goal checkers used to score
episode success \\
\midrule
\multicolumn{3}{@{}l}{\emph{Offline scripts (repository root)}} \\
\midrule
\texttt{run\_webarena.py} & 362 & Entry point: configuration resolution,
per-role LLM wiring, run orchestration \\
\texttt{build\_sft\_dataset.py} & 224 & Episodes $\rightarrow$ curriculum-ordered
SFT \texttt{jsonl} under prompt parity \\
\texttt{reannotate\_thoughts.py} & 186 & Rewrites each step's thought into
state / reflection / completion-check form \\
\texttt{judge\_episodes.py} & 108 & LLM re-verification of declared-success
episodes \\
\texttt{replay\_segmentation.py} & 131 & Offline replay of segmentation over
recorded episodes \\
\end{longtable}}

\section{Core data schema}
\label{app:design-schema}

Five dataclasses carry everything through the pipeline. Keeping them in one
module is what makes the LLM-induced artifacts inspectable, which in turn is
what makes the verification audit trail possible.

\begin{table}[H]
\centering
\caption[Core data schema]{Core data types and their fields.}
\label{tab:app-schema}
\small
\begin{tabular}{@{}L{3.1cm}L{3.6cm}L{6.8cm}@{}}
\toprule
\textbf{Type} & \textbf{Field} & \textbf{Meaning} \\
\midrule
\multirow{5}{*}{\texttt{Step}}
 & \texttt{action}         & Executed action string, e.g. \texttt{click('12')} \\
 & \texttt{url}            & URL \emph{after} the action \\
 & \texttt{obs\_summary}   & Short natural-language summary of the result \\
 & \texttt{dom\_signature} & Cheap structural fingerprint of the page \\
 & \texttt{axtree}         & Full accessibility tree (real backends only) \\
\midrule
\multirow{5}{*}{\texttt{SubTrajectory}}
 & \texttt{steps}        & Contiguous slice of a trajectory \\
 & \texttt{subgoal}      & Natural-language subgoal for the slice \\
 & \texttt{source\_site} & Domain the slice was collected from \\
 & \texttt{source\_task} & Task that produced the parent trajectory \\
 & \texttt{success}      & Only successes are eligible for discovery \\
\midrule
\multirow{4}{*}{\texttt{ParamSlot}}
 & \texttt{name}        & Lifted slot name, e.g. \texttt{username} \\
 & \texttt{slot\_type}  & \texttt{value} $\mid$ \texttt{selector} $\mid$ \texttt{url} \\
 & \texttt{description} & What the slot stands for \\
 & \texttt{example}     & A concrete observed value \\
\midrule
\multirow{7}{*}{\texttt{SkillCard}}
 & \texttt{name}, \texttt{description} & Human-readable identity \\
 & \texttt{precondition}       & Initiation set $I$, as checkable NL \\
 & \texttt{action\_template}   & Parameterized policy $\pi$ \\
 & \texttt{slots}              & Parameters of the template \\
 & \texttt{success\_predicate} & Termination test $\beta$ \\
 & \texttt{stats}              & Members, verification attempts/successes, sites \\
 & \texttt{verified}           & Admission flag \\
\midrule
\multirow{3}{*}{\texttt{Task}}
 & \texttt{instruction}  & Natural-language goal given to the solver \\
 & \texttt{skill\_chain} & Composed skill ids, in order \\
 & \texttt{difficulty}   & Chain length by default \\
\bottomrule
\end{tabular}
\end{table}

\noindent
The correspondence to the options framework is direct and intentional: a
\texttt{SkillCard} is an option $\langle I, \pi, \beta \rangle$ whose three
parts are the \texttt{precondition}, the \texttt{action\_template} together
with its \texttt{slots}, and the \texttt{success\_predicate}. The
\texttt{action\_signature} of a \texttt{SubTrajectory} --- the sequence of
action verbs alone --- is what makes clustering site-agnostic: logging in on
two different sites yields the same signature even though every selector and
value differs.

\section{Backend seam}
\label{app:design-backends}

Three interfaces separate the pipeline from anything that costs money, needs a
GPU, or requires a live site:

\begin{itemize}
  \item \texttt{LLM.complete\_json(system, user)} --- used by segmentation,
        abstraction, grounding, operator induction and both agents.
        \texttt{MockLLM} returns scripted responses; \texttt{OpenAILLM} calls
        any OpenAI-compatible endpoint, retries transient failures (rate
        limit, connection, timeout, 5xx --- never authentication or
        bad-request errors) up to three times with exponential backoff, and
        optionally attaches a screenshot for vision-capable models.
  \item \texttt{Embedder.encode(texts)} --- used by clustering and the
        predicate registry. \texttt{HashingEmbedder} is deterministic;
        \texttt{STEmbedder} wraps a sentence-transformers model.
  \item \texttt{WebEnv} with \texttt{reset}, \texttt{step} and \texttt{check}
        --- the pipeline needs nothing more from an environment. The mock
        provides scripted sites; \texttt{WebArenaBrowser} drives a real
        browser through BrowserGym.
\end{itemize}

\noindent
Because the seam is narrow, the entire discovery pipeline is exercised by
deterministic tests that assert exact skill counts and composition edges
without any API key, GPU or browser.

\section{Per-role model configuration}
\label{app:design-config}

Every LLM-backed helper is configured independently, so cheap work can go to a
cheap model while induction and judging use a stronger one. A role inherits any
field it does not override from the shared block; a global command-line flag
overrides all roles at once. Table~\ref{tab:app-roles} records the assignment
used for the runs in Chapter~\ref{ch:experiments}.

\begin{table}[H]
\centering
\caption[Per-role model assignment]{Per-role model assignment. ``Calls'' is the
order of magnitude of invocations per node, which is what justifies the
cost/strength trade-off in each row.}
\label{tab:app-roles}
\footnotesize
\begin{tabular}{@{}L{3.6cm}L{2.8cm}L{2.6cm}L{4.4cm}@{}}
\toprule
\textbf{Role} & \textbf{Model} & \textbf{Calls} & \textbf{Notes} \\
\midrule
\texttt{page\_explorer} & GPT-4o & once per node & Vision + Set-of-Marks,
\texttt{detail=low}, temperature 0.7 for proposal diversity \\
\texttt{solver} & GPT-4o-mini & up to
\texttt{max\_steps} per episode & Vision + Set-of-Marks, temperature 0.0 \\
\texttt{trajectory\_segmenter} & DeepSeek & per trajectory & Text only \\
\texttt{skill\_abstractor} & DeepSeek & per mature cluster & Text only \\
\texttt{skill\_verifier} & DeepSeek & per candidate skill & Text only \\
\texttt{operator\_inducer} & DeepSeek & per admitted skill & Text only \\
\texttt{episode\_judge} & DeepSeek & offline, per success & Text only \\
\texttt{thought\_reannotator} & DeepSeek & offline, per step & Text only \\
\bottomrule
\end{tabular}
\end{table}

\section{Run parameters}
\label{app:design-params}

Table~\ref{tab:app-params} lists the non-secret run parameters. Secrets ---
API keys, account passwords and site URLs --- live outside the repository in
an environment file. Precedence is command-line flag over configuration file
over built-in default.

\begin{table}[H]
\centering
\caption[Run parameters]{Run parameters and the values used.}
\label{tab:app-params}
\small
\begin{tabular}{@{}L{4.4cm}L{1.4cm}L{7.2cm}@{}}
\toprule
\textbf{Parameter} & \textbf{Value} & \textbf{Meaning} \\
\midrule
\texttt{max\_nodes}     & 500 & Node budget for one run \\
\texttt{max\_depth}     & 300 & Maximum frontier depth \\
\texttt{max\_steps}     & 15  & Solver action budget per episode \\
\texttt{max\_tasks}     & 9   & Tasks proposed per node \\
\texttt{maturity}       & 2   & Recurrences before a cluster becomes a skill \\
\texttt{task\_retries}  & 0   & Retries of a failed task (multiplies cost) \\
\texttt{goal\_chains\_per\_node} & 3 & New operator chains attempted per node \\
\texttt{predicate\_sim} & 0.6 & Cosine threshold of the predicate registry \\
\texttt{mine\_hrefs}    & true & Mine \texttt{<a href>} links into the frontier \\
\texttt{per\_template\_cap} & 3 & Nodes admitted per URL template \\
\texttt{max\_hrefs\_per\_page} & 100 & Links mined per page \\
\texttt{denylist} & --- & Blocks logout, pagination and sort URLs \\
\bottomrule
\end{tabular}
\end{table}

\chapter{Prompts}
\label{app:prompts}

Every prompt is reproduced verbatim, in the form the model actually receives.
They are given here rather than in Chapter~\ref{ch:method} because their exact
wording is the result of the behavioural fixes discussed in
Chapter~\ref{ch:discussion} --- the structured task classification, the
explicit answer-leaking prohibition and the duplicate-proposal ban in
Section~\ref{app:prompt-explorer} each replaced a formulation that a real run
showed to fail.

\section{Action space}
\label{app:prompt-actions}

Shared by the solver prompt and the offline dataset builder, so that a training
example never shows an action vocabulary different from the one available at
inference time.

\begin{lstlisting}[style=mystyle,numbers=none,basicstyle=\scriptsize\ttfamily]
Allowed actions (return exactly one): click('bid'), fill('bid','text'),
select_option('bid','opt'), press('bid','key_comb') -- e.g. press('12','Enter')
to submit a search box or form when there is no separate submit button,
goto('url'), go_back(), new_tab(), tab_focus(index), tab_close(),
send_msg_to_user('done message'), report_infeasible('reason') --
report_infeasible ONLY when the task cannot be done on this site/page.
\end{lstlisting}

\section{Page explorer}
\label{app:prompt-explorer}

\begin{lstlisting}[style=mystyle,numbers=none,basicstyle=\scriptsize\ttfamily]
PROPOSE concrete, feasible tasks a typical user could do on THIS page.
Return JSON {"tasks": [{"text": "...", "type": "INFO"|"NAV"|"MOD"}, ...]}.
Rules:
- Be SPECIFIC using names/values visible on the page (e.g. 'Add the Sony PS5
  to the cart', 'Open the post titled "Best budget GPU"'), NOT generic
  templates like 'add item to cart' or 'search for a product'.
- "type" must be exactly one of: INFO (read/extract info), NAV (go to a
  specific page), MOD (change content/state). List them INFO first, then NAV,
  then MOD.
- INFO tasks must name WHAT to find (e.g. 'Find the number of submissions for
  r/relationship_advice'), but must NEVER state the value/answer itself in the
  task text (NOT 'Find the number of submissions for r/relationship_advice
  (5,718 submissions)'). If the answer is already written in the task, it isn't
  a task -- whoever solves it must still have to read the page to answer it.
- Do NOT propose a NAV task to a page the user is already on (check the URL
  given below) -- e.g. if URL already ends in /wiki, do not propose 'Go to the
  Wiki page'.
- Do NOT re-propose a task you (or an earlier page) already proposed this run
  -- see 'ALREADY PROPOSED' below if present. Sitewide nav-menu links
  (header/footer, present on every page) are especially prone to this: once a
  destination has been proposed from one page, do not propose 'go to <that
  same destination>' again from a different page. Prefer content unique to
  THIS page.
- Prefer tasks that make progress or reveal new pages over pure
  search/sort/filter.
- Focus on common things a real user of THIS site would do; skip niche tasks.
- NEVER propose a task that would delete, deactivate, or close the user's
  account (e.g. 'Delete your account', 'Deactivate your profile'). This
  account is shared across exploration runs.
\end{lstlisting}

\noindent
In read-only mode the following note is appended, and any task the parser
cannot classify is dropped along with every \texttt{MOD} task:

\begin{lstlisting}[style=mystyle,numbers=none,basicstyle=\scriptsize\ttfamily]
READ-ONLY MODE: propose ONLY tasks with type INFO or NAV. Do NOT propose any
type MOD / content-modification tasks (no posting, commenting, voting, adding
to cart, or checkout).
\end{lstlisting}

\section{Solver}
\label{app:prompt-solver}

This is the prompt the fine-tuned model is trained against; the dataset builder
imports the constant rather than restating it, which is what guarantees prompt
parity between training and inference.

\begin{lstlisting}[style=mystyle,numbers=none,basicstyle=\scriptsize\ttfamily]
You are a web agent. Reach the GOAL by choosing the single best next action
from the page accessibility tree (elements shown as [bid] role 'name').
<ACTION SPACE, see B.1>
Think step by step; reflect on your past actions and any error from your last
action before deciding. Output ONE action as JSON
{"thought": "...", "action": "<one action>"}.
Examples: {"thought": "Open the product", "action": "click('12')"} ;
{"thought": "Goal met", "action": "send_msg_to_user('done: price is $19')"}.
Call send_msg_to_user ONLY when the goal is fully achieved.
For 'go to page X' goals: once the URL/page has changed to the requested
destination the goal is DONE -- call send_msg_to_user immediately; do NOT
click the same link again.
\end{lstlisting}

\section{Operator induction}
\label{app:prompt-operator}

The instruction to reuse a known predicate's exact text is the vocabulary-aware
mechanism of Section~\ref{sec:method-semantic}: embedding similarity alone
could not unify synonyms reliably, so the prompt is given the existing
vocabulary and asked to reuse it verbatim.

\begin{lstlisting}[style=mystyle,numbers=none,basicstyle=\scriptsize\ttfamily]
You convert a reusable web skill into a PLANNING OPERATOR. Given the skill and
the page state before/after it ran on a few sites, state the world conditions
in short, site-agnostic natural language. Preconditions = what must already be
true to apply the skill; effects = what becomes true after. Keep each predicate
a terse phrase (e.g. 'search results are shown', 'an item is in the cart').
Mark terminal=true only if an effect is a final/transactional state (order
placed, purchase confirmed). A list of KNOWN predicates from previously induced
operators may be provided. If a precondition or effect you want to state is
semantically equivalent to a known predicate, you MUST reuse that predicate's
exact text verbatim instead of rephrasing it. Only coin a new phrase when no
known predicate matches.
Return JSON {"preconditions": ["..."], "effects": ["..."], "terminal": bool}.
\end{lstlisting}

\noindent
A separate grounding check re-tests each induced predicate against the real
page states it was induced from, which is what makes the operator-quality
audit of Chapter~\ref{ch:discussion} possible:

\begin{lstlisting}[style=mystyle,numbers=none,basicstyle=\scriptsize\ttfamily]
You verify PLANNING PREDICATES against real page states. You are given a list
of predicates (each tagged precondition or effect) and, for each occurrence of
the skill, the page's accessibility tree BEFORE and AFTER the skill ran. For a
precondition, check it against the BEFORE tree of every occurrence; for an
effect, check it against the AFTER tree of every occurrence.
Return JSON {"results": [{"predicate": "<exact text>", "type":
"precondition"|"effect", "member_holds": [true|false, ...]}, ...]}, one entry
per predicate, with member_holds having exactly one boolean per occurrence in
the order given.
\end{lstlisting}

\section{Episode judge}
\label{app:prompt-judge}

The judge is deliberately told to grade the final state rather than the agent's
own claim --- the failure mode it exists to catch.

\begin{lstlisting}[style=mystyle,numbers=none,basicstyle=\scriptsize\ttfamily]
You are a strict grader of a web agent's completed task. You are given the GOAL
the agent was told to accomplish, the actions it took before declaring the task
done, and the FINAL page (URL + accessibility tree) where it declared done.
Decide whether the goal was genuinely achieved by the final state -- not
whether the agent SAID it was done. Common false positives: the agent gave up
and declared done anyway; it answered an INFO question with a guess not
actually shown on the final page; it navigated to the wrong page.
Return JSON {"verdict": "PASS"|"FAIL", "reason": "<one sentence>"}.
\end{lstlisting}

\section{Thought re-annotator}
\label{app:prompt-reannotator}

The subgoal enters this prompt but never the training prompt: at inference time
no component supplies ``the current subgoal'', so letting it into the input
would create exactly the train/serve mismatch the curation stage exists to
prevent.

\begin{lstlisting}[style=mystyle,numbers=none,basicstyle=\scriptsize\ttfamily]
You REWRITE one step's internal reasoning for a web agent's training data. You
are given the step's GOAL, its SUBGOAL (or, for the final step of a task, a
note that a final decision is being made), the action PRIOR to this one and its
outcome (if any), the current page, and the action this step took. Produce a
terse one-sentence STATE assessment of what the page currently shows. If the
PRIOR action's outcome shows no progress or an error, write a one-sentence
REFLECTION on why (empty string if the prior action clearly succeeded, or this
is the first step). If this step is the FINAL step of the task (marked below),
also judge in COMPLETION_CHECK whether the goal is actually satisfied by the
current page -- not just whether the agent believes so (empty string
otherwise).
Return JSON {"state": "<one short sentence>", "reflection": "<sentence or empty
string>", "completion_check": "<sentence or empty string>"}.
\end{lstlisting}

\chapter{Detailed Run Statistics}
\label{app:runs}

This appendix records the full statistics of the exploration runs summarized in
Section~\ref{sec:exp-data}. All twelve runs are reported, including the short
and interrupted ones, so that the collection effort is visible in full rather
than through its best runs only.

\section{Exploration runs}
\label{app:runs-explore}

Runs 1--9 explored the Reddit instance; runs 10--12 explored GitLab. Runs 4--6
are task-generation passes that reuse an existing graph and therefore have no
node counts of their own. A status of \emph{interrupted} means the run was
stopped manually once it had produced enough trajectories, not that it failed;
run~10 is the single run that ended on an error, a browser navigation timeout
that the node-level error handling converted into a skipped node before the
run was stopped.

\begin{table}[H]
\centering
\caption[Per-run exploration statistics]{Per-run exploration statistics.
``Tasks'' counts each task once, by final outcome.}
\label{tab:app-runs}
\small
\begin{tabular}{@{}lrrrrrrrr@{}}
\toprule
& & & \multicolumn{3}{c}{\textbf{Tasks}} & & & \\
\cmidrule(lr){4-6}
\textbf{Run} & \textbf{Nodes} & \textbf{Eps.} & \textbf{ok} & \textbf{fail} &
\textbf{inf.} & \textbf{Skills} & \textbf{Ops} & \textbf{Site} \\
\midrule
1  &   1 &   6 &   3 &   3 &  0 &  1 &  1 & Reddit \\
2  &   1 &   8 &   7 &   1 &  0 &  1 &  1 & Reddit \\
3  &  38 & 108 &  71 &  37 &  0 &  5 &  4 & Reddit \\
7  & 128 & 514 & 316 & 198 &  0 & 13 & 13 & Reddit \\
8  &  11 &  43 &  23 &  20 &  0 &  1 &  1 & Reddit \\
9  &  62 & 228 & 182 &  18 & 28 &  4 &  4 & Reddit \\
10 &   7 &  31 &  29 &   1 &  1 &  3 &  3 & GitLab \\
11 &  27 & 134 & 109 &  22 & 15 &  5 &  5 & GitLab \\
12 &  26 & 102 &  87 &  30 &  4 &  4 &  3 & GitLab \\
\midrule
\textbf{Total} & \textbf{301} & \textbf{1174} & \textbf{827} & \textbf{330} &
\textbf{48} & \textbf{37} & \textbf{35} & \\
\bottomrule
\end{tabular}
\end{table}

\noindent
Two observations are worth recording. First, the largest run (run~7, 128 nodes)
predates the infeasibility outcome, which is why its failure count is
proportionally the highest: episodes that a later run would have classified as
\emph{infeasible} are counted as failures there. Second, node duplication is
real but modest --- 29 nodes across all runs were skipped because their initial
accessibility tree hashed identically to an already-visited node, confirming
that distinct URLs rendering identical content is a genuine phenomenon on these
sites but not a dominant cost.

\section{Task generation passes}
\label{app:runs-gen}

Runs 4--6 measure task proposal and verification in isolation, over an already
explored graph. Table~\ref{tab:app-gen} breaks their outcomes down by task type.
The pattern is consistent across all three passes: \texttt{INFO} tasks verify
most reliably, \texttt{MOD} tasks least, and roughly one proposal in ten is
skipped as stale because the page it referred to had changed since the graph
was recorded.

\begin{table}[H]
\centering
\caption[Task generation and verification by type]{Task generation and
verification by task type, runs 4--6.}
\label{tab:app-gen}
\small
\begin{tabular}{@{}lrrrrrrrr@{}}
\toprule
& & \multicolumn{2}{c}{\textbf{INFO}} & \multicolumn{2}{c}{\textbf{NAV}} &
\multicolumn{2}{c}{\textbf{MOD}} & \\
\cmidrule(lr){3-4}\cmidrule(lr){5-6}\cmidrule(lr){7-8}
\textbf{Run} & \textbf{Att.} & \textbf{ok} & \textbf{fail} & \textbf{ok} &
\textbf{fail} & \textbf{ok} & \textbf{fail} & \textbf{Stale} \\
\midrule
4 & 92 & 43 &  9 & 12 & 8 & 16 & 4 &  8 \\
5 & 89 & 45 &  4 & 18 & 2 & 13 & 7 & 11 \\
6 & 85 & 37 &  8 & 13 & 7 & 15 & 5 & 15 \\
\midrule
\textbf{Total} & \textbf{266} & \textbf{125} & \textbf{21} & \textbf{43} &
\textbf{17} & \textbf{44} & \textbf{16} & \textbf{34} \\
\bottomrule
\end{tabular}
\end{table}

\noindent
Aggregated, \texttt{INFO} tasks verify at 86\%, \texttt{NAV} at 72\% and
\texttt{MOD} at 73\%. That \texttt{NAV} is no easier than \texttt{MOD} is the
counter-intuitive entry: navigation tasks fail not because the destination is
hard to reach but because the agent frequently cannot tell that it has already
arrived, which is precisely the behaviour the explicit ``once the URL has
changed the goal is DONE'' clause in the solver prompt
(Section~\ref{app:prompt-solver}) was added to correct.

\section{Composed goal tasks}
\label{app:runs-goal}

No composed goal task was attempted in any of the twelve runs: every recorded
episode, in every run directory, carries the low-level task tag, and the
composed-goal counters of all twelve run summaries are zero. The cause is
structural rather
than a defect in the chaining procedure. An edge in the operator graph requires
one operator's effect predicate to coincide with another's precondition, and
the predicate registry canonicalizes predicates by embedding similarity alone,
which --- as Section~\ref{sec:disc-operators} reports --- cannot reliably
separate true synonyms from unrelated phrases at any single threshold. In
practice almost every induced operator received a private predicate
vocabulary. Run~7, the largest graph with thirteen operators, contains exactly
one effect$\rightarrow$precondition edge, and that edge only came into
existence when the thirteenth and final operator was induced, by which point
the run had been stopped.

The consequence for the dataset is direct: the trajectories used for
fine-tuning are drawn entirely from low-level, page-scoped tasks, and the
long-horizon composed tasks that the operator graph is designed to produce are
not represented in it. The chaining mechanism itself is implemented and
covered by deterministic tests on the mock backend, where it produces the
expected chains and their sub-chains; what has not yet been demonstrated is
that a live site yields an operator graph connected enough for it to fire.
Section~\ref{sec:disc-operators} discusses the predicate-unification work this
implies.

\section{Per-task evaluation results}
\label{app:runs-eval}

\textit{[To be added once the evaluation sweeps for the remaining sites and the
statistical repeats of Section~\ref{sec:exp-stats} are complete. This section
will list, for each of the 180 official tasks per site, the baseline and
fine-tuned outcome, so that the aggregate success rates in
Section~\ref{sec:exp-main} can be audited task by task.]}

\chapter{Reproducibility Guide}
\label{app:repro}

The pipeline is reproducible end to end from a clean checkout. This appendix
gives the exact command sequence; the repository accompanying this thesis
contains the code, the configuration file used for the reported runs, and the
tutorial that this guide condenses.

\section{Environment}
\label{app:repro-env}

The mock pipeline needs only \texttt{numpy} and runs without an API key, a GPU
or a browser --- it is the recommended way to inspect the pipeline's behaviour
before spending anything:

\begin{lstlisting}[style=mystyle,numbers=none,basicstyle=\scriptsize\ttfamily]
pip install numpy
python demo.py             # discover -> cluster -> abstract -> verify -> transfer
python explore.py          # exploration loop: frontier -> discover -> compose
python compare_three_way.py  # Flat vs AWM vs HSGE on identical input
python -m pytest hsge/tests/ -q
\end{lstlisting}

\noindent
The real pipeline additionally needs \texttt{sentence-transformers},
\texttt{openai}, \texttt{browsergym}, \texttt{pyyaml},
\texttt{python-dotenv} and \texttt{requests}, a self-hosted WebArena
instance, and API credentials placed in an environment file that is never
committed.

\section{Collecting trajectories}
\label{app:repro-collect}

\begin{lstlisting}[style=mystyle,numbers=none,basicstyle=\scriptsize\ttfamily]
# Wiring check: scripted LLM and fake browser, no cost, no live site.
python run_webarena.py --dry-run --out /tmp/wa_dry --max-nodes 2

# Smoke test against a live shared instance without writing to it.
python run_webarena.py --read-only --max-nodes 2 --max-steps 4

# Full collection run (all parameters from config.yaml).
python run_webarena.py
\end{lstlisting}

\noindent
Login for the four account-gated sites happens once and is cached as a browser
storage state; the cache is a live session credential and is excluded from
version control. Run a single worker: WebArena sites are shared mutable state,
and concurrent runs interfere with each other. Artifacts --- skills, operators,
graph, predicates, per-node task outcomes and a summary --- are written after
every node, so an interrupted run still leaves a consistent snapshot, and every
event is mirrored to a JSON-lines log.

\section{Curating the dataset}
\label{app:repro-curate}

Both quality passes mutate the recorded episodes in place, additively; run them
against a copy if the original artifacts must stay pristine.

\begin{lstlisting}[style=mystyle,numbers=none,basicstyle=\scriptsize\ttfamily]
python judge_episodes.py outputs/<run>          # re-check declared successes
python reannotate_thoughts.py outputs/<run>     # thought -> thought_v2
python build_sft_dataset.py outputs/<run>/episodes train.jsonl --require-judge
\end{lstlisting}

\noindent
\texttt{--require-judge} drops any declared-success episode that was never
judged or that the judge failed, while keeping episodes the agent reported as
infeasible --- judging does not apply to them and they are valid training
targets.

\section{Fine-tuning, serving and evaluation}
\label{app:repro-train}

Fine-tuning uses LoRA on the base model, with the split taken over whole
trajectories rather than individual steps so that no trajectory contributes
steps to both sides. The chat template must have thinking mode disabled: the
collected data contains no reasoning traces in the template's expected form,
and leaving it enabled trains the model to emit a section it has never seen a
target for. After merging the adapter, the model is served locally and
evaluated with the official benchmark harness rather than a reimplementation,
so that the reported numbers are comparable with published results.

\section{Reported artifacts}
\label{app:repro-artifacts}

The following are the concrete artifacts behind Chapter~\ref{ch:experiments}:
the twelve run directories of Appendix~\ref{app:runs}, each with its skills,
operators, graph, predicate registry, per-node outcomes and event log; the
recorded episodes; the assembled SFT dataset; the LoRA adapter and merged
model; and the evaluation outputs of the official harness.
